\section{总结与延伸}

本讲建立了 Transformer 的最小完整图景：

\begin{enumerate}
  \item \textbf{历史不是断点}：RNN/seq2seq 暴露 fixed-vector bottleneck，早期 attention 提供按需读取，self-attention 将其升级为主要序列层；
  \item \textbf{Attention 是 routing}：query 表达需求，key 提供匹配索引，value 提供传递内容；
  \item \textbf{设计空间早已存在}：soft/hard 与 global/local 分别控制可微选择和访问范围；
  \item \textbf{Self-attention 缩短路径}：token 可直接交互并行训练，但 full score matrix 带来二次成本；
  \item \textbf{完整 block 不只有 attention}：position、mask、residual、normalization 与 FFN 分别解决顺序、可见性和深层优化；
  \item \textbf{家族由信息边界定义}：encoder-only 双向表示，decoder-only causal 生成，encoder-decoder 处理条件 seq2seq；
  \item \textbf{GPT-3 改变任务接口}：decoder-only pretraining 让任务可通过 prompt 和 example 表达；
  \item \textbf{BERT 强化表示接口}：encoder-only MLM 学习双向 contextual representation；
  \item \textbf{跨领域不是复制}：每个领域都要重新定义 token、position、mask、objective 与 evaluation；
  \item \textbf{未来问题不由架构自动解决}：memory、efficiency、alignment 与 deployment reliability 仍需独立设计。
\end{enumerate}

\begin{importantbox}{手算与设计作业}
给定三个 token 的 $Q,K,V$，手算 score、scaling、softmax 与 output；分别加入 padding mask 和 causal mask，解释哪几个权重必须为零。随后选择图像、RL trajectory 或蛋白质序列之一，定义 token、position、attention pattern、objective 与 metric，并说明为何使用 encoder-only、decoder-only 或 encoder-decoder。
\end{importantbox}

\begin{warningbox}{最终自检}
若你只能复述“Transformer 用 attention”，说明还没有掌握接口。应能解释 basic attention 与 self-attention 的来源差异、mask 如何定义信息边界、GPT 与 BERT 为何共享 block 却服务不同任务，以及 $O(n^2)$ 是数学结构还是具体实现瓶颈。
\end{warningbox}

\subsection{拓展阅读}

{\small
\begin{itemize}[nosep,leftmargin=*]
  \item \href{https://arxiv.org/abs/1409.0473}{\textbf{Bahdanau et al. (2014)}}：encoder-decoder additive attention 与 alignment。
  \item \href{https://arxiv.org/abs/1508.04025}{\textbf{Luong et al. (2015)}}：global/local attention 与 multiplicative scoring。
  \item \href{https://arxiv.org/abs/1706.03762}{\textbf{Vaswani et al. (2017)}}：scaled dot-product、multi-head 与原始 encoder-decoder Transformer。
  \item \href{https://arxiv.org/abs/2005.14165}{\textbf{Brown et al. (2020)}}：GPT-3 与 few-shot/in-context evaluation。
  \item \href{https://arxiv.org/abs/1810.04805}{\textbf{Devlin et al. (2018)}}：BERT、MLM、NSP 与 fine-tuning interface。
  \item \href{https://web.stanford.edu/class/cs25/past/cs25-v1/}{\textbf{Stanford CS25 V1}}：原始 speaker series 与课程安排。
\end{itemize}
}

\end{document}
